\section{VALIDATION EVIDENCE, RECONSTRUCTABILITY, AND BENCHMARK READINESS}
\label{sec:validation_reproducibility}

Section~\ref{sec:metrics} showed that a reported relationship becomes useful
for comparison when its task, measurement plane, operating point, and
validation setting align. Validation now asks how far that operating point was
tested. Reconstructability asks how completely another group can recover the
configuration and evaluation path that produced it. This section examines the
settings and methods used across the corpus, field evidence reported for both
functional domains, and the artifact support available for reconstruction and
benchmark comparison.

\subsection{Validation Settings and Methods}

The corpus has a broad experimental base. A total of 156 studies reached at
least a laboratory experiment or proof of concept. Within the exclusive
maximum setting distribution, basic simulation or numerical evidence was the
highest setting in 32 studies and enhanced simulation or evaluation supported
by a dataset was highest in 18. Laboratory experiments or proofs of concept
accounted for 78 studies, controlled prototypes for 66, and field trials or
deployments for 12. Each study enters one category, so the distribution covers
all 206 studies.

These settings expose different parts of an O-ISAC system. Simulation studies
isolate waveform and device sensitivities under controlled conditions
\cite{OISAC_SCR00007,OISAC_SCR00274}. Laboratory and prototype systems reveal
conversion, synchronization, and receiver effects in physical implementations
\cite{OISAC_SCR00001,OISAC_SCR00210,OISAC_SCR00362}. Field studies extend the
evidence to access networks, deployed fiber, submarine monitoring, and outdoor
wireless links
\cite{OISAC_SCR00496,OISAC_SCR00631,OISAC_SCR00957,OISAC_SCR01054,OISAC_SCR01086}.

The highest setting describes the reach of the evidence. Its scientific
meaning still depends on the claim and the test design. A controlled prototype
may explore a wider parameter range than a field trial, while a laboratory
experiment may reveal communication and sensing interactions more directly
than a field dataset collected for one outcome. The setting therefore gains
meaning when it is read together with the tested conditions, available
controls, and functional coverage.

The overlapping method map provides a complementary view. Analytical treatment
appeared in 131 studies, numerical analysis in 14, simulation in 104, and
evaluation with a dataset in 13. Laboratory experiments appeared in 148
studies, while 83 used a prototype or testbed and 12 reported a field
experiment. Mixed validation appeared in 33 studies. These counts permit more
than one method per study because analysis, simulation, measurement, and
prototype work can contribute to the same evidence chain.

The two maps answer different questions. Laboratory work contributed to 148
studies, while laboratory experiment or proof of concept was the strongest
setting for 78. Prototype or testbed work contributed to 83 studies, while a
controlled prototype was the strongest setting for 66.

Each method serves a different analytical purpose. Analysis clarifies
assumptions, simulation supports controlled parameter studies, and laboratory
work reveals the behavior of physical components. Prototypes bring subsystems
together, while field studies introduce at least one operational boundary.
Their combined value grows when the configuration and conditions remain
traceable as the evidence moves from one method to another.

Section~\ref{sec:optical_modality_families} shows that realistic validation
conditions change with the platform. Fiber studies depend on topology, traffic,
launch and receiver conditions, and the disturbance context. VLC/LiFi studies
add layout, orientation, illumination, detector, and processing conditions.
Free space optical studies require geometry, pointing, target properties,
weather, and motion. Terahertz systems enabled by photonics must also trace
optical conversion, antenna alignment, and wireless propagation. A useful
validation record therefore identifies the environment, measurement plane,
target or disturbance, reference method, and exercised function.

\subsection{Field Evidence Across Communication and Sensing}

The highest validation setting and the coverage of both functions answer
related questions. For 12 studies, a field trial or deployment was the
strongest reported setting. Six of them also received TQAF validation score 3
because their records report field or deployment outcomes in each functional
domain. The assessment is applied to each study. Electronic Supplement
S-Appraisal reports the corresponding scores, and Electronic Supplement S7
identifies the 12 studies and the subset of six.

The six studies span terahertz and cellular links enabled by photonics, fiber
to the room access integrated with radio frequency identification, space laser
ranging and communication, multicore fiber deployed in the field, and
submarine cable monitoring
\cite{OISAC_SCR01054,OISAC_SCR00496,OISAC_SCR00631,OISAC_SCR00730,OISAC_SCR01086,OISAC_SCR00957}.
Their diversity shows that field evidence for both functional domains appears
across several platforms. Platform comparisons still follow the native
validation conditions described in Section~\ref{sec:optical_modality_families}.

Electronic Supplement S7 records relationship timing and separate source
locators for the two functions as unresolved. The six cases therefore
establish field or deployment outcomes in both functional domains at the study
level. A stronger claim of concurrent joint operation requires both outcomes
to be traceable to the same configuration, time interval, and condition set.

The distinction becomes important when evidence is produced at different
stages of a system. One fiber study combines field data with offline sensing
classification and illustrates how field collection and functional execution
can occupy different parts of the evaluation path \cite{OISAC_SCR01006}.
Segmented architectures add interfaces, timing, calibration, and conversion
conditions to that path. When these elements connect both outcomes to one
operating point, field coverage can support a stronger claim of joint
operation.

% ============================================================================
% PLANNED FIGURE 7 --- VALIDATION PROFILE
% Stable label fig:validation_profile
% Blueprint ID FIG-OISAC-06
% Placement here after settings, methods, and functional coverage.
% Data authority Phase F s5_validation_maturity.csv and
% s5_validation_types.csv with denominator 206.
%
% Purpose
% Separate the strongest setting reached from the overlapping methods used to
% build the evidence.
%
% Layout
% Use exactly two aligned panels. This is a profile and carries no quality
% sequence.
%
% Panel A shows the mutually exclusive maximum settings. Use 32 basic
% simulation or numerical, 18 enhanced simulation with dataset support, 78
% laboratory or proof of concept, 66 controlled prototype, and 12 field trial
% or deployment. State that 156 studies reached at least laboratory or proof
% of concept.
%
% Panel B shows the overlapping methods. Use analytical 131, numerical analysis
% 14, simulation 104, dataset evaluation 13, laboratory experiment 148,
% prototype or testbed 83, field experiment 12, and mixed validation 33. Counts
% remain separate and are not summed.
%
% Activation rule
% Once the figure is produced and audited, replace the full count inventories
% in the first and fourth paragraphs of subsection VI-A with the following
% lead in.
% Figure~\ref{fig:validation_profile} separates the strongest setting reached
% by each study from the methods used to build the evidence.
%
% Future caption
% Validation profile of the 206 studies. Panel A assigns each study to its
% strongest setting. Panel B shows the validation methods reported across the
% corpus and permits more than one method per study.
%
% Keep functional coverage in prose and Electronic Supplement S7. Keep
% artifact access in Table VI and TQAF profiles in Figure 4.
% Use direct labels and a grayscale safe design.
% ============================================================================

\subsection{Artifact Access and Reconstructability}

Validation describes what a study tested. Reconstructability describes how
completely another group can recover the operating point and the path from
input to reported result. Access to data, code, or models supports this task,
while physical experiments also depend on the system configuration,
measurement conditions, processing steps, and provenance.

Table~\ref{tab:artifact_reconstruction} summarizes the access routes reported
for data and for code or models in all 206 studies.

\begin{table}[!t]
\caption{Reported access to data and to code or models in the 206 studies}
\label{tab:artifact_reconstruction}
\centering
\footnotesize
\setlength{\tabcolsep}{3.2pt}
\renewcommand{\arraystretch}{1.04}
\begin{tabularx}{\columnwidth}{@{}
>{\raggedright\arraybackslash}X
>{\centering\arraybackslash}p{0.28\columnwidth}@{}}
\toprule
Reported access & Studies, $n$ (\%) \\
\midrule
\multicolumn{2}{@{}l}{\textit{Data}} \\
Unavailable or not reported & 145 (70.4) \\
Available on request & 41 (19.9) \\
Open & 13 (6.3) \\
Not applicable & 7 (3.4) \\
\addlinespace[2pt]
\multicolumn{2}{@{}l}{\textit{Code or model}} \\
Unavailable or not reported & 197 (95.6) \\
Available on request & 7 (3.4) \\
Partial components & 1 (0.5) \\
Not applicable & 1 (0.5) \\
\bottomrule
\end{tabularx}
\vspace{1pt}
\parbox{\columnwidth}{\raggedright\footnotesize\emph{Note.} The first row in
each group combines records coded as unavailable with records that did not
report an access route at extraction.\par}
\end{table}

The reported routes support different forms of reconstruction. Open data
appear across simulation, laboratory, prototype, computer vision, and field
data settings
\cite{OISAC_SCR00274,OISAC_SCR00362,OISAC_SCR00833,OISAC_SCR01006,OISAC_SCR01008}.
They support reanalysis of the released data, while complete system
reconstruction also requires configuration and execution details. Software
preserves the transformation from inputs to reported outputs, and one computer
vision study reported partial code or model components
\cite{OISAC_SCR00833}.

The reconstruction framework developed in this review adds the information
needed to recreate a physical experiment. It covers the system and hardware,
the scene and calibration, the waveform and processing, the metrics and data,
and the versioned release. These requirements extend the comparison profile in
Section~\ref{sec:foundations_comparison}. Electronic Supplements S-Data
Dictionary and S-Evidence provide the field definitions and record provenance
used in the review.

The TQAF reproducibility assessment reads artifact access together with
parameter reporting. It places 4 studies in the low category, 199 in the
adequate category, and 3 in the strong category. Electronic Supplement
S-Appraisal provides the study level scores. Thirteen studies report open data,
while three reach strong reproducibility. The difference reflects the
additional reconstruction information captured by the appraisal.

\subsection{Benchmark Readiness}

A reconstructable result becomes a reusable benchmark when another group can
apply the same task, condition set, baseline, and metric definition. The TQAF
benchmark assessment combines these requirements with artifact access,
validation, reproducibility, and comparison admissibility. It places 48
studies in the low category and 158 in the adequate category. The strong
category contains zero studies under the review criteria.

Under the review criteria, strong readiness depends on bringing an external or
common baseline, an open artifact, admissible outcomes in both functional
domains, and strong validation and reproducibility support into the same
evidence chain. Electronic Supplement S-Appraisal provides the study level
scores. The zero count directs attention from isolated demonstrations toward
shared tasks, baselines, artifacts, and evaluation conditions that other groups
can reuse.

These distributions describe validation and reporting patterns within the 206
included studies. The survey design supports conclusions about evidence reach,
functional coverage, artifact access, and reconstructability. Questions about
causal relations among artifact access, validation setting, and reported
performance require a different study design. Section~\ref{sec:discussion_roadmap}
turns the observed gaps into specific benchmark and reporting actions.

Overall, the corpus combines a broad laboratory and prototype base with a
smaller set of field and deployment results. Stronger cumulative evidence
emerges when both functions are tied to one traceable configuration and the
evaluation path is preserved through reusable artifacts. Section~\ref{sec:technologies_applications_6g}
builds on this validation map by examining the technologies and applications
supported by the reported evidence and their connection to 6G network roles.
